\documentclass[11pt]{article}
\usepackage{fontspec,amsmath,amssymb,amsthm,hyperref}
\usepackage[margin=1in]{geometry}
\setlength{\emergencystretch}{3em}
\title{The joint initial-state and proposal density}
\author{LangevinLean mathematical exposition}
\date{}
\begin{document}\maketitle
\noindent Source module: \texttt{UnadjustedLangevin/GaussianProposalPairDensity.lean}. Compile this proof document with XeLaTeX or LuaLaTeX. The accompanying HTML pairs the mathematical argument with the exact Lean statements and records the scope of the formalization.


For a measurable center \(c(x)\), independent \(X\sim\mu\) and \(Z\sim\mathcal N(0,I)\) produce the joint law \(\mu(dx)q_h(c(x),y)dy\).


\section{Disintegrate the Gaussian change of variables}


For \(X\sim\mu\) and independent \(Z\sim\mathcal N(0,I)\), fix \(x\) and apply the proved proposal-density identity to \(c(x)+\sqrt{2h}Z\). It gives conditional integration against \(q_h(c(x),y)dy\).



To obtain equality of joint measures, evaluate a measurable set \(A\subset H\times H\). Tonelli writes both sides as integrals over \(x\) of their \(y\)-sections. The fixed-center density theorem identifies each section, hence the measures agree. Measurability of \(c\) supplies measurability of the pair update and joint density; no probability or moment condition on \(\mu\) is added.



\noindent Lean declarations: \texttt{gaussian\_proposal\_pair\_density}.

\end{document}
